% SGA 3, 第 VI_B 篇, 第 1 轮英文重构.
% 范围: 第 5 节标题至注 5.4.1, 含底本编者注 33--36 和脚注中的定理 5.3A.
% 底本: Polo--Gille Exp6B-13oct24.pdf, 分篇第 25--26 页
% (合订本第 366--367 页; 原版页边印刷页码 351--352, 重版本页眉 355--356).
% 推论 5.5 始于分篇第 27 页 / 合订本第 368 页.
% OCR 与 Reinhold 英文 Markdown 仅为比对材料.
% 英文底本旧图注: 命题 5.1 图为底本 180 dpi 精确裁切:
% 分篇第 25 页, 1488x2105 渲染中的 x=565,y=1080,w=365,h=180.
% 本译稿保留下方已有的原生 TikZ 图码.

\section*{5. 群与齐性空间的分离性}
\phantomsection\label{sec:sga3-VIB-5}%
\addcontentsline{toc}{section}{5. 群与齐性空间的分离性}

\medskip
\phantomsection\label{prop:sga3-VIB-5.1}%
\noindent\textbf{命题 5.1.---}
\emph{\(S\)-群 \(G\) 分离的充要条件是 \(G\) 的单位截面为闭浸入.}

此条件是必要的 (EGA~I, 5.4.6); 由下列笛卡儿图可知它也是充分的
(参见 VI\(_A\), 0.3):
\begin{center}
  \phantomsection\label{diag:sga3-VIB-5.1-cartesian}%
  \begin{tikzcd}[row sep=large,column sep=huge]
    G\times_S G
      \arrow[r,"{\mu\circ(\operatorname{id}_G\times c)}"] &
    G \\
    G
      \arrow[u,"\Delta_{G/S}"]
      \arrow[r,"\pi"] &
    S
      \arrow[u,"\varepsilon"']
  \end{tikzcd}\,.
\end{center}

\medskip
\phantomsection\label{prop:sga3-VIB-5.2}%
\noindent\textbf{命题 5.2.---}
\emph{若 \(S\) 离散, 则每个 \(S\)-群都分离.}

事实上,
\[
  S=\coprod_{s\in S}\operatorname{Spec}\mathcal O_{S,s},
\]
且由 EGA~I, 5.5.5, 只须证明对每个 \(s\in S\),
\[
  G\times_S\operatorname{Spec}\mathcal O_{S,s}
\]
分离. 这由 VI\(_A\), \SGARefLink{sga3:VIA:0.3}{0.3} 得到,
因为 \(\mathcal O_{S,s}\) 是维数为 \(0\) 的局部环.%
\footnote{\label{note:sga3-VIB-33}编者注: 我们将``Artin 环''改为
``维数为 \(0\)'', 并在 VI\(_A\), \SGARefLink{sga3:VIA:0.3}{0.3} 中提及这一推广.}

\medskip
\phantomsection\label{thm:sga3-VIB-5.3}%
\noindent\textbf{定理 5.3.---}
\emph{设 \(S\) 是概形, \(G\) 是在 \(S\) 上局部有限呈示,
且在单位截面的某邻域中对 \(S\) 普遍开的 \(S\)-群概形,
\(X\) 是一个带有 \(G\) 作用的 \(S\)-概形, 使态射}
\[
  \Phi:G\times_S X\longrightarrow X\times_S X,
  \qquad (g,x)\longmapsto(gx,x),
\]
\emph{满射. 再假设对每个 \(s\in S\):}
\begin{enumerate}
\item[\textup{(i)}] \emph{存在 \(X\) 的一个在 \(S\) 上分离的开子概形
  \(U\), 使 \(U_s\) 在 \(X_s\) 中稠密;}
\item[\textup{\SGAInlineRefTarget{thm:sga3-VIB-5.3:item:ii}(ii)}]
  \emph{纤维 \(X_s\) 在 \(\kappa(s)\) 上局部有限型.}%
\footnote{\label{note:sga3-VIB-34}编者注: 我们通过假设 \(G\)
在单位截面的某邻域中对 \(S\) 普遍开, 并增加假设
\textup{\SGARefLink{thm:sga3-VIB-5.3:item:ii}{(ii)}}, 改正了原文;
见下文 O.\ Gabber 的例子, 它们表明没有附加假设时,
原文的陈述 \hyperref[thm:sga3-VIB-5.3]{5.3} 和 \hyperref[cor:sga3-VIB-5.4]{5.4} 不成立.
还指出, \hyperref[thm:sga3-VIB-5.3]{定理 5.3}
是下面\hyperref[thm:sga3-VIB-5.3A]{定理 5.3A} 的改写版本,
后者见于 SGAD 的 1965 年版, 归功于 M.\ Raynaud;
比较 \SGARefLink{sga3:x:lemma:8-5}{第~X 篇, 8.5} 和 8.8 中的注 \((*)\).

\smallskip
\phantomsection\label{thm:sga3-VIB-5.3A}%
\noindent\textbf{定理 5.3A (Raynaud).---}
\emph{设 \(G\) 是局部有限呈示、在 \(S\) 上普遍开且纤维连通的 \(S\)-群.
则 \(G\) 在 \(S\) 上分离. 更一般地, 每个在 \(S\) 上局部有限呈示,
并带有 \(G\) 的作用, 使态射
\(\Phi:G\times_SX\to X\times_SX\),
\((g,x)\mapsto(gx,x)\), 满射的 \(S\)-概形 \(X\), 都在 \(S\) 上分离.}}
\end{enumerate}
\emph{则 \(X\) 在 \(S\) 上分离.}

\medskip
\phantomsection\label{cor:sga3-VIB-5.4}%
\noindent\textbf{推论 5.4.---}
\emph{设 \(S,G,X\) 如 \SGARefLink{thm:sga3-VIB-5.3}{5.3} 的预备假设中所述.
再假设 \(X\) 的纤维局部有限型且连通. 则:}
\begin{enumerate}
\item[\textup{\SGAInlineRefTarget{cor:sga3-VIB-5.4:item:i}(i)}]
  \emph{\(X\) 在 \(S\) 上分离.}
\item[\textup{(ii)}] \emph{若 \(X\) 有一个在 \(S\) 上拟紧的开集
  \(V\), 它与每个非空纤维 \(X_s\) 相交,}%
\footnote{\label{note:sga3-VIB-35}编者注: 若没有此假设,
则有 O.\ Gabber 指出的如下反例.
设 \(G=S\) 是以 \(s\) 为闭点的局部概形,
使 \(S^\ast=S-\{s\}\) 不拟紧,
并设 \(X\) 是 \(\{s\}\) 与 \(S^\ast\) 的不交并.}
\emph{则 \(X\) 在 \(S\) 上拟紧.}
\end{enumerate}

\emph{证明.}
\textup{\SGARefLink{cor:sga3-VIB-5.4:item:i}{(i)}}
设 \(s\in S\) 使 \(X_s\ne\varnothing\).
由 \(\Phi\) 作基变换所得的态射
\[
  \Phi_s:
  G_s\times_{\kappa(s)}X_s
  \longrightarrow
  X_s\times_{\kappa(s)}X_s
\]
满射. 由于 \(X_s\) 连通, 由 VI\(_A\),
\SGARefLink{prop:sga3-VIA-2.5.4}{2.5.4}, 它不可约.
因此, 若 \(U\) 是 \(X\) 的一个使 \(U_s\) 非空的仿射开集,
则 \(U_s\) 在 \(X_s\) 中稠密, 定理即可应用.

为证明 \textup{\SGARefLink{concl:sga3-VIA-2.5.4-ii}{(ii)}},
可假设 \(S\) 仿射. 则 \(V\) 拟紧,
且由 \hyperref[prop:sga3-VIB-3.5]{3.5},
\(G\) 有一个包含 \(\underline{G^0}\) 的拟紧开集 \(U\).
设 \(s\in S\) 使 \(X_s\ne\varnothing\).
则 \(X_s\) 不可约 (VI\(_A\), \SGARefLink{prop:sga3-VIA-2.6.6}{2.6.6}),
并且由于 \(U_s\) 包含 \(\underline{G_s^0}\), 态射
\[
  U_s\times_{\kappa(s)}V_s\longrightarrow X_s,
  \qquad (g,x)\longmapsto gx,
\]
满射 (VI\(_A\), \SGARefLink{prop:sga3-VIA-2.6.4}{2.6.4}). 因而态射
\[
  U\times_SV\longrightarrow X
\]
满射. 由于 \(U\) 和 \(V\) 拟紧, \(X\) 也拟紧.
由于已假设 \(S\) 仿射, 因而分离,
可知 \(X\) 在 \(S\) 上拟紧 (参见 EGA~I, 6.6.4\textup{(v)}).

\medskip
\phantomsection\label{rem:sga3-VIB-5.4.1}%
\noindent\textbf{注 5.4.1.---}%
\footnote{\label{note:sga3-VIB-36}编者注: 此注为我们新增;
比较编者\hyperref[note:sga3-VIB-34]{注 34}.}
在证明过程中我们将看到, 只假设下述条件,
\hyperref[thm:sga3-VIB-5.3]{定理 5.3} 的结论仍成立:

\begin{quote}
\textup{(i')} 存在 \(X\) 的一个在 \(S\) 上分离的开子概形 \(U\),
使 \(U_s\) 在 \(X_s\) 的每个不可约分支中稠密.
\end{quote}

若 \(X_s\) 局部地只有有限个不可约分支,
如假设 \textup{(ii)} 下的情形, 则这由
\textup{\SGARefLink{cor:sga3-VIB-5.4:item:i}{(i)}} 得到.
后面我们还将看到, 在每个纤维 \(X_s\) 都拟分离且连通的假设下,
\SGARefLink{cor:sga3-VIB-5.4}{5.4} 仍成立.
